\documentclass[11pt]{article}
\usepackage[T1]{fontenc}
\usepackage{lmodern}
\usepackage{amsmath,amssymb,amsthm,mathtools,bm}
\usepackage{enumitem}
\newtheorem{theorem}{Theorem}[section]
\newtheorem{lemma}[theorem]{Lemma}
\theoremstyle{definition}
\newtheorem{definition}[theorem]{Definition}
\newtheorem{example}[theorem]{Example}
\newcommand{\R}{\mathbb{R}}
\newcommand{\N}{\mathbb{N}}
\newcommand{\abs}[1]{\left\lvert #1 \right\rvert}
\newcommand{\norm}[1]{\left\lVert #1 \right\rVert}
\DeclareMathOperator{\Var}{Var}
\DeclareMathOperator*{\argmax}{arg\,max}
\begin{document}
\section{Sequences}
\begin{definition}
A sequence $(a_n)_{n\in\N}$ in $\R$ \emph{converges} to $L$ if for every $\varepsilon>0$ there is $N$ such that $\abs{a_n-L}<\varepsilon$ for all $n\geq N$.
\end{definition}

\begin{theorem}[Uniqueness of limits]\label{thm:unique}
A convergent sequence has exactly one limit.
\end{theorem}
\begin{proof}
Suppose $a_n\to L$ and $a_n\to L'$. Then
\begin{align}
  \abs{L-L'} &\leq \abs{L-a_n}+\abs{a_n-L'} \label{eq:tri}\\
             &< 2\varepsilon
\end{align}
for all large $n$, so $L=L'$ by \eqref{eq:tri}.
\end{proof}

\begin{example}
The sequence $a_n = \dfrac{1}{n}$ converges to $0$, while $b_n = (-1)^n$ does not converge.
\end{example}

The variance of $X$ is $\Var(X) = \mathbb{E}\big[(X-\mu)^2\big]$ and the maximiser is
\[
  x^\star = \argmax_{x\in[0,1]} \bigl( f(x) - \tfrac{1}{2}\norm{x}^2 \bigr),
  \qquad
  f(x) = \begin{cases} x^2 & x \geq 0,\\ -x & x < 0. \end{cases}
\]

\begin{enumerate}[label=\roman*)]
  \item Show that $\sum_{k=1}^{n} k = \frac{n(n+1)}{2}$ for all $n \in \N$.
  \item Compute $\lim_{x\to 0} \frac{\sin x}{x}$ and $\int_0^1 e^{-x^2}\,\mathrm{d}x$ numerically.
  \item Let $\bm{v} \in \R^3$; find $\norm{\bm{v}}$ when $\bm{v} = (1,2,2)^{\top}$.
\end{enumerate}

\begin{lemma}
For matrices $A = \begin{pmatrix} a & b \\ c & d \end{pmatrix}$ we have $\det A = ad - bc$ and, if $\det A \neq 0$, $A^{-1} = \frac{1}{\det A}\begin{pmatrix} d & -b \\ -c & a\end{pmatrix}$.
\end{lemma}
By Theorem~\ref{thm:unique} the limit is unique. Note that $\overbrace{1+1+\cdots+1}^{n\text{ times}} = n$ and $\underbrace{x \cdot x \cdots x}_{k} = x^k$. The construction follows \cite{rudin1976}; see also \cite[Ch.~2]{abbott2015}.

\begin{thebibliography}{9}
\bibitem{rudin1976} W.~Rudin, \emph{Principles of Mathematical Analysis}, 3rd ed., McGraw--Hill, 1976.
\bibitem{abbott2015} S.~Abbott, \emph{Understanding Analysis}, 2nd ed., Springer, 2015.
\end{thebibliography}
\end{document}
